\chapter{Существующие компоненты}

\begin{warnbox}[Глава заменена роадмапом экосистемы]
Состояние на 29 сентября 2026 года. Эта глава заменена роадмапом экосистемы (вики, \texttt{roadmap/ecosystem-roadmap.ru.md}, черновик v4: порядок приоритетов, модуль хоста сессий SH, этапы M0--M8, треки PC, KP и CM) и репозиториями модулей; карта репозиториев по плоскостям дана в \texttt{README.md} вики. С тех пор, как глава написана, хостинг сессий выносится в модуль agent-session-host, процесс описывает конфигурация процесса (curator-playbook), доверие вынесено в curator-trust, главный цикл узла в curator-keeper, сообщения и аренды между оркестраторами в waggle, а собственный агентский цикл и минимальный runtime на нём отложены (curator-keeper, §7.7). Текст главы сохранён как история; где он расходится с роадмапом или спецификацией модуля, правы они.
\end{warnbox}

Протокол не строится с нуля. Четыре компонента уже существуют, и большая часть механики протокола есть их композиция. Эта глава описывает каждый в той мере, в какой он нужен для протокола, и даёт карту соответствий.

\section{Карта соответствий}

\begin{center}
\small
\begin{tabular}{p{4.2cm} p{4.4cm} p{4.6cm}}
\toprule
Понятие протокола & Компонент & Что берётся \\
\midrule
запись, подпись, каноническая форма & реестр аудита Curator & CCJ-1, конверт Ed25519, ротация ключей \\
цепочка, дерево, снапшот, защита от отката & реестр аудита Curator & лог прозрачности, Merkle, снапшоты, высшая принятая версия \\
план сессии, digest, только сужение & curator-agent-runtime & скомпилированный план, эпоха плана, делегирование вниз \\
receipt, грант как потолок & curator-agent-runtime & execution-receipt, call plan, host assurance \\
аренда, fencing, чекпоинт & ax & lease record, эпоха аренды, checkpoint record \\
передача сессии, конвертация & ax & handoff, force takeover, clone между агентами \\
задача, оценка, класс, ревью & борд задач & README и progress, конфигурация спавна \\
контракт задачи, версии & борд задач, реестр & README с критериями как тело контракта; версии как записи журнала \\
ключ идентичности, ротация & реестр аудита Curator & набор пинов, ротация через перекрытие, отзыв удалением \\
скилл remote-shell, пины & Curator & closure, артефакты с хешами, провайдеры \\
\bottomrule
\end{tabular}
\end{center}

\section{Curator: скиллы, closure и реестр}

Curator есть протокол пакетов для агентских скиллов: единица распространения содержит инструкции для модели, провайдеров инструментов и артефакты, всё с идентичностью по содержимому. Компилятор разрешает closure: полный набор пинов, из которого собирается план сессии. Два свойства важны для протокола.

Первое: runtime не поставляет инструментов, инструменты приходят как скиллы. Скилл удалённого shell есть обычный пакет с провайдерами (локальный, по SSH, в контейнере) и артефактами (бинарь демона под каждую платформу с хешами). Замена провайдера или демона это изменение хеша в closure, видимое в диффе плана. Это ровно та форма, в которой публичный клиент получает две свои привязки.

Второе: реестр аудита. Записи с подписями Ed25519 поверх канонического JSON, идентификатор ключа как префикс хеша ключа, ротация ключей через перекрытие, лог прозрачности с непрерывными номерами и хешем предыдущей записи, дерево Меркла, подписанные снапшоты с корнем, длиной и головой, защита от отката через высшую принятую версию как долговременное состояние, федерация реестров с правилом deny-wins, офлайн-бандлы с контрподписью при импорте. Это готовый журнал уровня L0. Ему не хватает четырёх вещей, и все они добавляются, а не переписываются: типы записей для задач, сертификат кворума как порог эндорсментов (формат уже хранит эндорсменты с их конвертами подписей), доказательства включения для светлых клиентов, секвенсер как роль сервиса и сертификат остановки как запись с особым правилом принятия.

Инвентарь ограничений Curator честно помечает на macOS ограничение файловой системы и сети как недоступное в непривилегированном портативном режиме. Демон среды поэтому есть verified host execution provider в терминах Curator: отдельно установленный доверенный компонент с точным хешем бинаря, permit и receipt. Seatbelt живёт в нём, а не в портативном режиме.

\section{Runtime агента}

Runtime разделён на три части: компилятор, который строит план сессии заранее и офлайн; брокер, который владеет консолью и границей исполнения; ядро, которое обрабатывает вывод модели и не имеет ни ключей, ни файловой системы. Четыре свойства runtime протокол берёт как есть.

План сессии скомпилирован, неизменяем в течение своей эпохи (это первое из трёх значений слова эпоха) и имеет digest; каждый запрос, вызов и результат несёт идентификатор эпохи и digest плана, и результаты из прошлой эпохи не принимаются. Протокол связывает digest плана с записью назначения, и среда сверяет его на каждом вызове.

Делегирование только сужает: план дочерней сессии есть подмножество родительского, ни одного инструмента, гранта или лиза сверх. Протокол поднимает это правило на уровень задач.

Грант есть потолок, план вызова есть применённый профиль, receipt есть наблюдённое. Для shell точность всегда неизвестна, поэтому shell работает на потолке; отсюда требование, чтобы потолок был санбоксом задачи, а не машиной. Схема receipt уже содержит digest плана, digest плана вызова, идентичность провайдера, отчёт об ограничении, дайджесты входа и выхода.

Доверенный путь: компонент, обрабатывающий вывод модели, не отвечает на одобрения. В протоколе тот же принцип: оркестратор, будучи моделью, не подписывает кворум своих записей.

Host assurance сообщается как полное, частичное или отсутствующее, и вызывающий, требующий полного, отказывается. Это инвариант I9 в готовом виде.

\section{ax: сессии, аренды, передача}

ax есть менеджер сессий агентов: у логической сессии ровно один активный владелец, реплики только читают, владение оформлено арендой с эпохой (второе значение слова) и случайным fencing-токеном, события владельца несут эпоху, пиры отвергают события младшей эпохи, но сохраняют их в ветке расхождений. Плановая передача требует доказательства покоя: ввод заблокирован, foreground и background idle, ноль открытых дочерних процессов и хендлов, стабильный дайджест хранилища. Принудительный захват есть отдельное явное действие. Чекпоинты адресуются по содержимому и служат валидированной базой передачи.

Клон между агентами (например, из одной агентской среды в другую) создаёт новую логическую сессию и не переносит аренду, владение бордом, одобрения, учётные данные и полномочия незавершённых операций. Это утилита конвертации сессий, о которой говорит протокол, и она уже воплощает правило: с сессией не путешествует авторитет. Общий принцип ax, транскрипты и записи есть content authority, а индекс есть перестраиваемая проекция, есть то же самое, что журнал истина, сессия кэш.

Удалённая точка входа ax это RPC по стандартным потокам поверх SSH без открытых портов; демон среды следует тому же образцу.

\section{Борд задач}

Борд живёт в git-репозитории проекта как дерево EPIC, STORY, TASK. У задачи есть README с описанием, областью, пришпиленной к коммиту, и критериями приёмки, и progress со статусом, флагом обязательности ревью, классом, оценкой в единицах сложности, блокировками и чеклистом; ресурсы задачи хранят логи запусков, вердикты ревью и патчи. Конфигурация спавна задаёт максимальную параллельность, изоляцию в копиях рабочего дерева с командами проверки, базовую ветку интеграции и потолки моделей по ролям разработчика и ревьювера.

Это почти содержимое журнала минус подписи, кворум и репликация. Оценка задачи есть цена, флаг ревью есть класс, потолки моделей по ролям есть принимаемые модели манифеста, команды проверки есть проверки из профиля контракта. Не хватает одного: версионирования критериев. README задачи сегодня правится как обычный файл; в журнале его изменение должно становиться новой версией контракта с отдельной авторизацией. Материализация журнала как борда в git, описанная в главе 5, есть просто признание того, что борд уже есть.

\section{Что берётся как есть и что расширяется}

Как есть: подписи, каноническая форма, цепочка, снапшоты, защита от отката; план сессии и его digest; делегирование только вниз; receipts; аренды с fencing; чекпоинты; клон без переноса полномочий; формат борда.

Расширяется: реестр до журнала с типами записей, кворумом, доказательствами включения и секвенсером; спавн борда до оркестратора с планировщиком по ёмкости и областям; скилл shell до восстановленной сессии с подписью на машине владельца.

Новое: демон среды, секвенсер журнала, протокол ноды. Следующая глава описывает их.
